\begin{theorem}[MetaCIC parallel-diamond seed closure review]
\label{thm:metacic-parallel-diamond-frontier-seed-closure-review}
Every seed-level read of an accepted $\MetaCICParallelDiamondFrontierUp$ packet must cite the critical-pair envelope, residual-join locality, candidate-normalization scope, finite-residual witness, retained-premise boundary, and closurestatus child routes before the packet can be consumed as a parallel-diamond frontier seed surface. The read is the finite route through $P,K,J,R,S,C,B,O,H,T,G,N$; it is not residual-substitution discharge, subject reduction, unrestricted confluence, evaluator equality, quotient term equality, global normalization, or kernel consistency.
\end{theorem}
\begin{proof}
Begin with the critical-pair and candidate-normalization rows. The critical-pair envelope and candidate-normalization scope keep $K,J,R,S,C,B,O$ attached to the retained premise row $P$ before any frontier consumer reads the candidate join.

Next read the residual and retained-premise routes. Residual-join locality, finite-residual witness consumption, and retained-premise boundary theorems keep the local join, residual witness, obstruction, and premise rows in the same finite packet.

Finally use the closurestatus child routes. They force every seed-status read through the displayed carrier, obligation, closurestatus export, and guard children, so the structural rows $H,T,G,N$ preserve only the listed seed surface and cannot promote it to any refused global MetaCIC endpoint.
\end{proof}
